% alg-p3-05: 因式分解决策树（流程图）
\documentclass[tikz,border=4pt]{standalone}
\usepackage{ctex}
\usepackage{amsmath}
\usetikzlibrary{arrows.meta, shapes.geometric}
\begin{document}
\begin{tikzpicture}[scale=1.0, thick,
  box/.style={rectangle, draw, thick, rounded corners=3pt,
              minimum width=3.0cm, minimum height=0.65cm,
              font=\footnotesize, align=center},
  dia/.style={diamond, draw, thick, aspect=2.8,
              minimum width=3.0cm, minimum height=0.8cm,
              font=\footnotesize, align=center, fill=yellow!20},
  res/.style={rectangle, draw, thick, rounded corners=3pt,
              minimum width=2.8cm, minimum height=0.6cm,
              font=\footnotesize, align=center, fill=blue!10},
  arr/.style={-{Stealth[length=4pt]}, thick}
]

  % 顶部起点
  \node[box, fill=gray!20] (start) at (0,0) {拿到多项式};

  % 问1：公因式？
  \node[dia] (q1) at (0,-1.3) {有公因式？};
  \node[res, fill=orange!15] (yes1) at (3.5,-1.3)
    {提出公因式\\（再从头走）};

  % 问2：项数
  \node[dia] (q2) at (0,-2.8) {项数？};

  % 两项分支
  \node[dia] (q3) at (-3.5,-4.2) {两平方相减？};
  \node[res] (r3y) at (-3.5,-5.4) {平方差\\$a^2{-}b^2{=}(a{+}b)(a{-}b)$};
  \node[res, fill=gray!15] (r3n) at (-5.8,-4.2) {一般不可分};

  % 三项分支
  \node[dia] (q4) at (0,-4.2) {中间项是\\$2$倍乘积？};
  \node[res] (r4y) at (0,-5.4) {完全平方\\$(a{\pm}b)^2$};
  \node[res] (r4n) at (2.4,-5.4) {十字相乘法};

  % 四项分支
  \node[dia] (q5) at (3.8,-4.2) {能二二分组？};
  \node[res] (r5y) at (3.8,-5.4) {分组提公因式};
  \node[res, fill=gray!15] (r5n) at (6.2,-4.2) {调整顺序};

  % 终止
  \node[box, fill=green!18] (done) at (0,-6.8) {每个因式不可再分 $\to$ 完成};

  % 关键提醒
  \node[font=\footnotesize, red, align=center] at (0,-7.55)
    {每完成一步都问：还能继续分解吗？};

  % ---- 箭头 ----
  \draw[arr] (start) -- (q1);
  \draw[arr] (q1) -- node[left, font=\tiny] {无} (q2);
  \draw[arr] (q1) -- node[above, font=\tiny] {有} (yes1);
  % yes1 → q2（回溯）
  \draw[arr] (yes1.south) -- ++(0,-0.4) -| (q2.east);

  % q2 → 三分支
  \draw[arr] (q2) -- node[above left, font=\tiny] {两项} (q3);
  \draw[arr] (q2) -- node[right, font=\tiny] {三项} (q4);
  \draw[arr] (q2) -- node[above right, font=\tiny] {四项} (q5);

  % 两项分支
  \draw[arr] (q3) -- node[left, font=\tiny] {是} (r3y);
  \draw[arr] (q3) -- node[above, font=\tiny] {否} (r3n);

  % 三项分支
  \draw[arr] (q4) -- node[left, font=\tiny] {是} (r4y);
  \draw[arr] (q4) -- node[above, font=\tiny] {否} (r4n);

  % 四项分支
  \draw[arr] (q5) -- node[left, font=\tiny] {是} (r5y);
  \draw[arr] (q5) -- node[above, font=\tiny] {否} (r5n);

  % → done
  \draw[arr] (r3y.south) |- (done.west);
  \draw[arr] (r4y) -- (done);
  \draw[arr] (r4n.south) |- (done);
  \draw[arr] (r5y.south) |- (done.east);

\end{tikzpicture}
\end{document}
